\documentclass[11pt]{article}

\usepackage[T1]{fontenc}
\usepackage{lmodern}
\usepackage{amsmath,amssymb,amsthm,mathtools}
\usepackage[margin=1in]{geometry}
\usepackage{microtype}
\usepackage[hidelinks]{hyperref}

\newtheorem{theorem}{Theorem}
\newtheorem{lemma}[theorem]{Lemma}
\newtheorem{proposition}[theorem]{Proposition}
\newtheorem{corollary}[theorem]{Corollary}
\theoremstyle{remark}
\newtheorem{remark}[theorem]{Remark}

\newcommand{\R}{\mathbb R}

\title{Purely Analytic Positivity of the Upward-Closed Scalar Gap}
\author{}
\date{Global middle-ratio analytic module}

\begin{document}
\maketitle

\begin{abstract}
The angular-block reduction for the middle-ratio spherical-shell proxy
leaves one smooth scalar inequality.  This note proves that inequality on
the upward-closed region
\(
 \rho_c\leq\rho\leq39/50,
 \ K\geq k_{11}(\rho)
\)
without interval arithmetic or a finite numerical certificate.  The loss
integral is first differentiated with respect to the spectral cutoff.  A
turning-point lower bound and an exact beta integral show that the scalar gap
is strictly increasing in that cutoff for every $K>12$.  At the base curve,
a rational secant majorant gives a uniform analytic estimate for the loss.
The remaining one-variable inequality is then reduced to positivity of a
degree-nine polynomial; all of its Bernstein coefficients on one rational
interval are displayed and are positive.  Thus there is no upper-frequency
cutoff in the scalar-positivity theorem.
\end{abstract}

\section{The scalar gap}

For \(0\leq s\leq1\), put
\begin{equation}
 G_1(s):=\frac1\pi\left(\sqrt{1-s^2}
                 -s\arccos s\right).
 \label{eq:G1}
\end{equation}
Fix
\begin{equation}
 \rho_c:=\frac1{1+2\pi},
 \qquad
 k_{11}(\rho):=
 \sqrt{\frac{\pi^2}{(1-\rho)^2}+132},
 \label{eq:constants}
\end{equation}
and consider
\begin{equation}
 \mathcal R_{\rm c}:=
 \left\{(\rho,K):
 \rho_c\leq\rho\leq\frac{39}{50},\quad
 K\geq k_{11}(\rho)\right\}.
 \label{eq:region}
\end{equation}
The scalar gap left by the angular-block estimate is
\begin{equation}
 \begin{aligned}
 \mathcal S(\rho,K):={}&
 \frac{(1+2\rho^2)K^2}{12}
 -\mathcal J(\rho,K)
 -\frac{\pi\rho K}{4(1-\rho)}+\frac12\\
 &+\frac{2\pi\rho K}{\arccos\rho}
 \left(
 \frac{KG_1(\rho)}{4(1+2KG_1(\rho))}
 +\frac3{32}
 \right),
 \end{aligned}
 \label{eq:S}
\end{equation}
where
\begin{equation}
 \mathcal J(\rho,K):=
 \frac{K^2}{2}\int_\rho^1
 \frac{s}{(1+2KG_1(s))^2}\,ds.
 \label{eq:J}
\end{equation}

\begin{theorem}[Upward-closed scalar positivity]
\label{thm:scalar-positive}
For every \(\rho_c\leq\rho\leq39/50\) and every
\(K\geq k_{11}(\rho)\),
\begin{equation}
 \boxed{\mathcal S(\rho,K)>0.}
 \label{eq:target}
\end{equation}
\end{theorem}

We use only the elementary rational estimates
\begin{equation}
 \frac{157}{50}<\pi<\frac{22}{7},
 \qquad
 \frac{140}{99}<\sqrt2,
 \qquad
 \frac53<\sqrt3.
 \label{eq:rational-constants}
\end{equation}
The two square-root estimates follow immediately by squaring; for the
closer one we record the exact margin
\[
 2\cdot99^2-140^2=2>0.
\]
The displayed bounds for \(\pi\) are the classical Archimedean bounds
(with the lower one weakened to \(157/50\)).

\section{Strict monotonicity in the spectral cutoff}

The identity
\begin{equation}
 G_1(s)=\frac1\pi\int_s^1\arccos u\,du
 \label{eq:G-integral}
\end{equation}
and \(1-\cos v\leq v^2/2\) imply
\begin{equation}
 G_1(s)\geq
 \frac{2\sqrt2}{3\pi}(1-s)^{3/2}.
 \label{eq:turning-lower}
\end{equation}
Indeed, \(\arccos u\geq\sqrt{2(1-u)}\), and integration gives
\eqref{eq:turning-lower}.

\begin{lemma}[Derivative of the boundary-layer loss]
\label{lem:J-derivative}
For every fixed \(0<\rho<1\) and every \(K>0\),
\begin{equation}
 \partial_K\mathcal J(\rho,K)
 =K\int_\rho^1
 \frac{s}{(1+2KG_1(s))^3}\,ds.
 \label{eq:J-derivative}
\end{equation}
Moreover, if \(K>12\), then
\begin{equation}
 \int_\rho^1
 \frac{s}{(1+2KG_1(s))^3}\,ds
 <\frac{88}{567}.
 \label{eq:cubic-loss}
\end{equation}
\end{lemma}

\begin{proof}
Differentiating \eqref{eq:J} under the integral sign gives
\[
 \begin{aligned}
 \partial_K\mathcal J
 &=
 K\int_\rho^1\frac{s}{(1+2KG_1(s))^2}\,ds
 -2K^2\int_\rho^1
 \frac{sG_1(s)}{(1+2KG_1(s))^3}\,ds\\
 &=K\int_\rho^1
 \frac{s}{(1+2KG_1(s))^3}\,ds,
 \end{aligned}
 \]
which proves \eqref{eq:J-derivative}.

Set \(a=4\sqrt2/(3\pi)\).  By \eqref{eq:turning-lower},
the change of variable \(t=1-s\), and \(s\leq1\),
\begin{align}
 \int_\rho^1
 \frac{s\,ds}{(1+2KG_1(s))^3}
 &\leq
 \int_0^\infty\frac{dt}{(1+aKt^{3/2})^3}\notag\\
 &=
 \frac{8\pi}{27\sqrt3}(aK)^{-2/3}.
 \label{eq:beta-bound}
\end{align}
The last equality is the beta integral
\[
 \frac23(aK)^{-2/3}
 B\left(\frac23,\frac73\right)
 =\frac{8\pi}{27\sqrt3}(aK)^{-2/3}.
 \]
At \(K=12\), the cube of the last member of
\eqref{eq:beta-bound} is
\[
 \frac{\pi^5}{59049\sqrt3}.
 \]
Using \eqref{eq:rational-constants},
\[
 \frac{\pi^5}{59049\sqrt3}
 <
 \frac{5153632}{1654060905}
 <
 \left(\frac{88}{567}\right)^3,
 \]
where the second strict inequality is the exact identity
\[
 \left(\frac{88}{567}\right)^3
 -\frac{5153632}{1654060905}
 =\frac{27812576}{44659644435}>0.
 \]
The right side of \eqref{eq:beta-bound} decreases with \(K\), proving
\eqref{eq:cubic-loss}.
\end{proof}

\begin{proposition}[Monotonicity of the scalar gap]
\label{prop:K-monotone}
For every \(\rho_c\leq\rho\leq39/50\) and every
\(K\geq k_{11}(\rho)\),
\begin{equation}
 \partial_K\mathcal S(\rho,K)>0.
 \label{eq:S-monotone}
\end{equation}
Consequently,
\begin{equation}
 \mathcal S(\rho,K)\geq
 \mathcal S\bigl(\rho,k_{11}(\rho)\bigr).
 \label{eq:base-reduction}
\end{equation}
\end{proposition}

\begin{proof}
Write
\[
 \theta:=\arccos\rho,\qquad y:=KG_1(\rho).
 \]
The derivative of the final line of \eqref{eq:S} is
\begin{equation}
 \frac{2\pi\rho}{\theta}
 \left(
 \frac{y(1+y)}{2(1+2y)^2}+\frac3{32}
 \right)
 \geq\frac{3\rho}{8},
 \label{eq:endpoint-derivative}
\end{equation}
because \(\theta\leq\pi/2\).  Also
\[
 -\frac{\pi\rho}{4(1-\rho)}
 >-\frac{11\rho}{14(1-\rho)}.
 \]
Every \(K\geq k_{11}(\rho)\) in the stated ratio interval is greater than
\(12\); this also
follows from the stronger estimate established in
Lemma~\ref{lem:uniform-loss} below.  Lemma~\ref{lem:J-derivative}
therefore gives
\begin{align}
 \partial_K\mathcal S
 &>
 K\left(\frac{1+2\rho^2}{6}-\frac{88}{567}\right)
 +\frac{3\rho}{8}
 -\frac{11\rho}{14(1-\rho)}
 \notag\\
 &=
 K\left(\frac{13}{1134}+\frac{\rho^2}{3}\right)
 +\frac{3\rho}{8}
 -\frac{11\rho}{14(1-\rho)}
 \notag\\
 &>
 \frac{26}{189}+4\rho^2+\frac{3\rho}{8}
 -\frac{11\rho}{14(1-\rho)}.
 \label{eq:S-derivative-lower}
\end{align}

It remains to check the sign of the last expression.  Multiplication by
\(1-\rho\) produces
\begin{equation}
 p(\rho):=
 (1-\rho)
 \left(\frac{26}{189}+4\rho^2+\frac{3\rho}{8}\right)
 -\frac{11\rho}{14}.
 \label{eq:p}
\end{equation}
The bound \(\pi<22/7\) gives
\(\rho_c>7/51\).  On the larger rational interval
\([7/51,39/50]\), set
\[
 u:=\frac{\rho-7/51}{39/50-7/51}\in[0,1].
 \]
In the degree-three Bernstein basis
\(B_{j,3}(u)=\binom3j u^j(1-u)^{3-j}\), exact expansion gives
\begin{equation}
 p(\rho)=\sum_{j=0}^3 b_jB_{j,3}(u),
 \label{eq:p-Bernstein}
\end{equation}
with
\begin{equation}
 \begin{array}{c|c}
 j&b_j\\ \hline
 0&\dfrac{335005}{2785671}\\[3pt]
 1&\dfrac{32947387}{196635600}\\[3pt]
 2&\dfrac{281783183}{578340000}\\[3pt]
 3&\dfrac{231517}{13500000}
 \end{array}
 \label{eq:p-coefficients}
\end{equation}
All four coefficients are positive.  Since the Bernstein basis functions
are nonnegative and sum to one, \(p(\rho)>0\).  Equation
\eqref{eq:S-derivative-lower} proves \eqref{eq:S-monotone}, and
\eqref{eq:base-reduction} follows by integration in \(K\).
\end{proof}

\section{An exact rational upper bound for the loss}

\begin{lemma}[Uniform loss bound]
\label{lem:uniform-loss}
For every \(\rho_c\leq\rho\leq39/50\) and every
\(K\geq k_{11}(\rho)\),
\begin{equation}
 K>\frac{241}{20}
 \label{eq:K-lower}
\end{equation}
and
\begin{equation}
 \int_\rho^1
 \frac{x}{(1+2KG_1(x))^2}\,dx
 <\frac{17}{100}.
 \label{eq:loss-17}
\end{equation}
\end{lemma}

\begin{proof}
Since \(\rho\geq\rho_c\),
\[
 \frac{\pi}{1-\rho}\geq
 \frac{\pi}{1-\rho_c}=\pi+\frac12>\frac{91}{25}.
 \]
Consequently
\[
 K^2\geq k_{11}(\rho)^2
 >132+\left(\frac{91}{25}\right)^2
 =\left(\frac{241}{20}\right)^2+\frac{471}{10000},
 \]
which proves \eqref{eq:K-lower}.  The rational bounds in
\eqref{eq:rational-constants} now give
\begin{equation}
 \frac{4\sqrt2K}{3\pi}
 >
 \frac{4(140/99)(241/20)}{3(22/7)}
 =\frac{23618}{3267}
 =\frac{36}{5}+\frac{478}{16335}.
 \label{eq:turning-coefficient}
\end{equation}

Put \(y=1-x\).  Equations \eqref{eq:turning-lower} and
\eqref{eq:turning-coefficient} yield
\begin{align}
 \int_\rho^1\frac{x\,dx}{(1+2KG_1(x))^2}
 &<
 \int_0^{1-\rho}
 \frac{1-y}{(1+(36/5)y^{3/2})^2}\,dy\notag\\
 &\leq
 \int_0^1
 \frac{1-y}{(1+(36/5)y^{3/2})^2}\,dy.
 \label{eq:universal-loss}
\end{align}
We bound the last integral by exact rational arithmetic.  Set
\[
 c:=\frac{36}{5},\qquad
 \phi(z):=\frac1{(1+z)^2},
 \]
and use the rational grid
\begin{equation}
 (t_i)_{i=0}^{10}=
 \left(
 0,\frac14,\frac5{16},\frac38,\frac7{16},\frac12,
 \frac9{16},\frac58,\frac34,\frac78,1
 \right).
 \label{eq:secant-grid}
\end{equation}
For \(z_i=ct_i^3\), let
\[
 \beta_i:=
 \frac{\phi(z_{i+1})-\phi(z_i)}{z_{i+1}-z_i},
 \qquad
 \alpha_i:=\phi(z_i)-\beta_i z_i.
 \]
Because \(\phi''(z)=6(1+z)^{-4}>0\), its secant satisfies
\[
 \phi(z)\leq\alpha_i+\beta_i z
 \qquad(z_i\leq z\leq z_{i+1}).
 \]
After \(y=t^2\), the integral in \eqref{eq:universal-loss} is
\[
 \int_0^1 2t(1-t^2)\phi(ct^3)\,dt.
 \]
Thus it is at most
\begin{equation}
 Q:=\sum_{i=0}^9
 \bigl(F_i(t_{i+1})-F_i(t_i)\bigr),
 \label{eq:Q-def}
\end{equation}
where the exact polynomial antiderivative is
\begin{equation}
 F_i(t):=
 \alpha_i\left(t^2-\frac{t^4}{2}\right)
 +2c\beta_i\left(\frac{t^5}{5}-\frac{t^7}{7}\right).
 \label{eq:Fi}
\end{equation}
Every number in \eqref{eq:Q-def} is rational.  Substitution of the grid
\eqref{eq:secant-grid} gives
\begin{equation}
 Q=
 \frac{
 1305767238061446154064110434835938779253983373178558326887
 }{
 7706446260271137290663804190970969288616921855152043745280
 },
 \label{eq:Q-value}
\end{equation}
and direct subtraction gives
\begin{equation}
 \frac{17}{100}-Q
 =
 \frac{
 21643130923235926743681388145629999054466710986445549053
 }{
 38532231301355686453319020954854846443084609275760218726400
 }>0.
 \label{eq:Q-gap}
\end{equation}
Equations \eqref{eq:universal-loss}--\eqref{eq:Q-gap} prove
\eqref{eq:loss-17}.
\end{proof}

\section{The base curve}

Fix \(\rho\) and now set
\begin{equation}
 t:=\sqrt{1-\rho},
 \qquad
 k:=k_{11}(\rho)
 =\sqrt{132+\frac{\pi^2}{t^4}},
 \qquad
 \rho=1-t^2.
 \label{eq:t-base}
\end{equation}
The middle-ratio range lies strictly inside the rational interval
\begin{equation}
 \frac7{15}<t<\frac{14}{15}.
 \label{eq:t-interval}
\end{equation}
Indeed,
\[
 t^2\geq\frac{11}{50},
 \qquad
 \frac{11}{50}-\left(\frac7{15}\right)^2
 =\frac1{450}>0,
\]
while \(\rho\geq\rho_c>7/51\) gives
\[
 t^2=1-\rho<\frac{44}{51}
 <\left(\frac{14}{15}\right)^2,
 \qquad
 \left(\frac{14}{15}\right)^2-\frac{44}{51}
 =\frac{32}{3825}>0.
\]

\begin{lemma}[Analytic endpoint payment]
\label{lem:endpoint-payment}
For every \(\rho_c\leq\rho\leq39/50\), let \(t\) and \(k\) be as in
\eqref{eq:t-base}, and let
\begin{equation}
 C(\rho):=\frac{100\rho^2-1}{600}
 \label{eq:C}
\end{equation}
and
\begin{equation}
 E(t):=
 \frac{11}{14t^2}
 -\frac{14}{15(1+2t)}
 -\frac{3}{8t}.
 \label{eq:E}
\end{equation}
Then \(C(\rho)>0\), \(E(t)>0\), and
\begin{equation}
 \mathcal S(\rho,k)>
 C(\rho)k^2-\rho E(t)k+\frac12.
 \label{eq:S-base-lower}
\end{equation}
\end{lemma}

\begin{proof}
First, \(\rho\geq\rho_c>7/51>1/10\), so \(C(\rho)>0\).
Lemma~\ref{lem:uniform-loss} and
\[
 \frac{1+2\rho^2}{12}-\frac{17}{200}
 =\frac{100\rho^2-1}{600}
 \]
show that the first two terms in \eqref{eq:S} are strictly greater
than \(C(\rho)k^2\).

Set \(\theta=\arccos\rho\) and \(Y=kG_1(\rho)\).
For \(0\leq t\leq1\),
\begin{equation}
 \theta=\arccos(1-t^2)
 =2\arcsin\frac{t}{\sqrt2}
 \leq\frac{\pi t}{2}.
 \label{eq:theta-bound}
\end{equation}
To see the last inequality, use convexity of \(\arcsin\) on
\([0,1/\sqrt2]\): its graph lies below the chord joining
\((0,0)\) to \((1/\sqrt2,\pi/4)\).
Furthermore, \eqref{eq:turning-lower} and
\(k>\pi/t^2\) give
\begin{equation}
 Y>
 \frac{2\sqrt2}{3}t
 >\frac{14}{15}t.
 \label{eq:Y-lower}
\end{equation}
Since \(Y\mapsto Y/(1+2Y)\) is increasing,
\eqref{eq:theta-bound}--\eqref{eq:Y-lower} imply
\begin{align}
 &\frac{2\pi\rho k}{\theta}
 \left(\frac{Y}{4(1+2Y)}+\frac3{32}\right)\notag\\
 &\qquad\quad>
 \rho k\left(
 \frac{14}{15(1+2t)}+\frac{3}{8t}
 \right).
 \label{eq:endpoint-lower}
\end{align}
Also
\[
 -\frac{\pi\rho k}{4t^2}>
 -\frac{11\rho k}{14t^2}.
 \]
Combining these inequalities with \eqref{eq:S} proves
\eqref{eq:S-base-lower}.

Finally,
\begin{equation}
 840t^2(1+2t)E(t)
 =660+1005t-1414t^2.
 \label{eq:E-numerator}
\end{equation}
The quadratic on the right is concave, so its minimum on
\([7/15,14/15]\) occurs at an endpoint.  Its endpoint values are
\[
 \frac{184739}{225},
 \qquad
 \frac{82406}{225},
 \]
respectively.  Both are positive, proving \(E(t)>0\).
\end{proof}

\begin{proposition}[Rational positivity on the base curve]
\label{prop:base-positive}
For every \(\rho_c\leq\rho\leq39/50\),
\begin{equation}
 \mathcal S\bigl(\rho,k_{11}(\rho)\bigr)>0.
 \label{eq:base-positive}
\end{equation}
\end{proposition}

\begin{proof}
The elementary inequality
\[
 \sqrt{x^2+132}<x+\frac{66}{x}\qquad(x>0)
 \]
and \eqref{eq:rational-constants} give
\begin{align}
 k^2&>
 132+\frac{(157/50)^2}{t^4},
 \label{eq:k-square-lower}\\
 k&<
 \frac{22}{7t^2}+\frac{3300t^2}{157}.
 \label{eq:k-upper}
\end{align}
Since \(C(\rho)>0\) and \(\rho E(t)>0\),
Lemma~\ref{lem:endpoint-payment} yields
\begin{equation}
 \mathcal S(\rho,k)>R(t),
 \label{eq:R-lower}
\end{equation}
where
\begin{equation}
 \begin{aligned}
 R(t):={}&
 C(1-t^2)
 \left(132+\frac{(157/50)^2}{t^4}\right)\\
 &-(1-t^2)E(t)
 \left(\frac{22}{7t^2}+\frac{3300t^2}{157}\right)
 +\frac12.
 \end{aligned}
 \label{eq:R}
\end{equation}

It remains only exact polynomial algebra.  Put
\[
 H(t):=t^4(1+2t)R(t).
 \]
Using \eqref{eq:E-numerator}, its factorized form is
\begin{equation}
 \begin{aligned}
 H(t)={}&
 C(1-t^2)(1+2t)
 \left(132t^4+\left(\frac{157}{50}\right)^2\right)\\
 &-\frac{(1-t^2)(660+1005t-1414t^2)}{840}
 \left(\frac{22}{7}+\frac{3300}{157}t^4\right)\\
 &+\frac12t^4(1+2t).
 \end{aligned}
 \label{eq:H-factorized}
\end{equation}
For an independently checkable expansion,
\begin{equation}
 \begin{aligned}
 H(t)={}&
 -\frac{20642567}{24500000}
 -\frac{6205067}{12250000}t
 +\frac{547983}{122500}t^2
 -\frac{1033727}{367500}t^3\\
 &+\frac{34911551}{16485000}t^4
 +\frac{187093801}{8242500}t^5
 +\frac{8679}{1099}t^6\\
 &-\frac{138149}{2198}t^7
 -\frac{2101}{157}t^8
 +44t^9.
 \end{aligned}
 \label{eq:H-expanded}
\end{equation}

Set
\[
 u:=\frac{t-7/15}{7/15}\in[0,1].
 \]
In the degree-nine Bernstein basis,
\begin{equation}
 H(t)=\sum_{j=0}^9 d_jB_{j,9}(u),
 \qquad
 B_{j,9}(u)=\binom9j u^j(1-u)^{9-j},
 \label{eq:H-Bernstein}
\end{equation}
where exact conversion of \eqref{eq:H-expanded} gives
\begin{equation}
 \begin{array}{c|c}
 j&d_j\\ \hline
 0&\dfrac{65408801807437}{9463832437500000}\\[3pt]
 1&\dfrac{299549381760793}{1135659892500000}\\[3pt]
 2&\dfrac{15904699805720773}{28391497312500000}\\[3pt]
 3&\dfrac{2765605256590021}{3154610812500000}\\[3pt]
 4&\dfrac{33203943051258761}{28391497312500000}\\[3pt]
 5&\dfrac{38636221171331747}{28391497312500000}\\[3pt]
 6&\dfrac{2538003206646073}{1892766487500000}\\[3pt]
 7&\dfrac{28758289199580211}{28391497312500000}\\[3pt]
 8&\dfrac{12920786461478803}{28391497312500000}\\[3pt]
 9&\dfrac{3428732855971679}{9463832437500000}
 \end{array}
 \label{eq:H-coefficients}
\end{equation}
Every \(d_j\) is positive.  The Bernstein basis functions are
nonnegative and sum to one; hence \(H(t)>0\) throughout
\([7/15,14/15]\).  Since \(t^4(1+2t)>0\), this proves \(R(t)>0\).
Equation \eqref{eq:R-lower} proves \eqref{eq:base-positive}.
\end{proof}

\begin{proof}[Proof of Theorem~\ref{thm:scalar-positive}]
Proposition~\ref{prop:K-monotone} reduces every point of
the upward-closed middle-ratio region to the base curve
\(K=k_{11}(\rho)\).
Proposition~\ref{prop:base-positive} proves strict positivity on that
curve.  Therefore \(\mathcal S(\rho,K)>0\) throughout
\(\mathcal R_{\rm c}\), with no upper restriction on \(K\).
\end{proof}

\begin{remark}[Nature of the certificate]
The rational secants and Bernstein expansions above are finite exact
algebra, not interval arithmetic: convexity proves the functional
majorant on each secant interval, and positivity follows from explicitly
displayed rational coefficients.  No floating-point evaluation, sampled
minimum, or numerical root isolation is used as a premise.
\end{remark}

\end{document}
